% ==========================================================================
%  Cabecalho LaTeX do curso "Trabalhando com Texto em Python I"
%  Injetado pelo pandoc (-H) ao gerar o .tex de cada unidade.
% ==========================================================================
\usepackage[T1]{fontenc}
\usepackage{lmodern}
\usepackage[scaled=0.92]{helvet}
\renewcommand{\familydefault}{\sfdefault}
\usepackage{microtype}
\usepackage{amssymb}
\usepackage{newunicodechar}
\usepackage{xcolor}
\usepackage[most]{tcolorbox}
\tcbuselibrary{skins,breakable}
\usepackage{titlesec}
\usepackage{enumitem}

% -------- paleta --------
\definecolor{accent}{HTML}{6C4CF1}
\definecolor{cyan}{HTML}{00B4D8}
\definecolor{subink}{HTML}{2B3242}
\definecolor{notec}{HTML}{3B82F6}
\definecolor{tipc}{HTML}{10B981}
\definecolor{warnc}{HTML}{F59E0B}

% -------- emojis / simbolos unicode que podem aparecer no texto --------
% simbolos com significado
\newunicodechar{✅}{{\color{tipc}\checkmark}}
\newunicodechar{✔}{{\color{tipc}\checkmark}}
\newunicodechar{☐}{$\square$}
\newunicodechar{κ}{\ensuremath{\kappa}}
\newunicodechar{≈}{\ensuremath{\approx}}
\newunicodechar{−}{\ensuremath{-}}
\newunicodechar{↓}{\ensuremath{\downarrow}}
\newunicodechar{→}{\ensuremath{\rightarrow}}
\newunicodechar{←}{\ensuremath{\leftarrow}}
\newunicodechar{↹}{\ensuremath{\hookrightarrow}}
\newunicodechar{°}{\textdegree}
\newunicodechar{·}{\textperiodcentered}
\newunicodechar{´}{'}
\newunicodechar{‑}{-}
% emojis decorativos -> removidos no PDF
\newunicodechar{💡}{}
\newunicodechar{🚀}{}
\newunicodechar{📝}{}
\newunicodechar{⚠}{}
\newunicodechar{🐍}{}
\newunicodechar{🎉}{}
\newunicodechar{🍕}{}
\newunicodechar{👍}{}
\newunicodechar{💋}{}
\newunicodechar{😚}{}
\newunicodechar{🚂}{}
\newunicodechar{❤}{}

% -------- caixas de callout (usadas pelo filtro callouts.lua) --------
\newtcolorbox{notebox}{enhanced,breakable,colback=notec!6,colframe=notec,
  boxrule=0pt,leftrule=4pt,arc=3pt,left=10pt,right=10pt,top=7pt,bottom=7pt,
  before skip=10pt,after skip=12pt}
\newtcolorbox{tipbox}{enhanced,breakable,colback=tipc!7,colframe=tipc,
  boxrule=0pt,leftrule=4pt,arc=3pt,left=10pt,right=10pt,top=7pt,bottom=7pt,
  before skip=10pt,after skip=12pt}
\newtcolorbox{warnbox}{enhanced,breakable,colback=warnc!8,colframe=warnc,
  boxrule=0pt,leftrule=4pt,arc=3pt,left=10pt,right=10pt,top=7pt,bottom=7pt,
  before skip=10pt,after skip=12pt}
\newtcolorbox{objbox}{enhanced,breakable,colback=accent!5,colframe=accent,
  boxrule=0pt,leftrule=4pt,arc=3pt,left=10pt,right=10pt,top=7pt,bottom=8pt,
  before skip=10pt,after skip=14pt}

% -------- titulos coloridos --------
\titleformat{\section}{\sffamily\Large\bfseries\color{accent}}{\thesection}{0.6em}{}
\titleformat{\subsection}{\sffamily\large\bfseries\color{subink}}{\thesubsection}{0.5em}{}
\titleformat{\subsubsection}{\sffamily\normalsize\bfseries\color{cyan}}{\thesubsubsection}{0.5em}{}
\titlespacing*{\section}{0pt}{2.2ex plus 1ex minus .2ex}{1.1ex}

% codigo inline colorido
\definecolor{inlink}{HTML}{5A34D6}
\let\oldtexttt\texttt
\renewcommand{\texttt}[1]{{\color{inlink}\oldtexttt{#1}}}
